\chapter{Larger Documents and Own Commands}\label{chap:larger}

A lab report fits into one file.
A thesis of fifty pages does not, or at least it should not:
the editor stays faster and clearer with several files,
and two people can work on different chapters without getting in each other's way.
This chapter shows how to split a document,
and how to write your own commands, which every larger document collects over time.

\section{Own commands}\label{sec:own-commands}

\code|\newcommand| defines a command of your own, in the preamble:

\owncodetitled{ch10-newcommand-preamble.tex}{Preamble}
\owncodetitled{ch10-newcommand.tex}{Document}
\ownrender{ch10-newcommand.pdf}

\begin{codeblock}[style=shell]
\newcommand{\name}[number of arguments]{definition}
\end{codeblock}

\begin{itemize}
    \item The name must not exist yet; otherwise the error is \texttt{Command \textbackslash name already defined}.
        \code|\renewcommand| replaces an existing command instead, as in \autoref{sec:emph-preamble}.
    \item In the definition, \code|#1| stands for the first argument, \code|#2| for the second, up to \code|#9|.
    \item With a default value, the first argument becomes optional (see below):
        \code|\vessel| prints \enquote{the mug} and \code|\vessel[flask]| \enquote{the flask}.
    \item Like every command without an argument, \code|\Tenv| swallows the space after it in the text
        (\autoref{sec:commands}); in math mode, that does not matter.
\end{itemize}

\begin{codeblocktitled}{Preamble}
\newcommand{\vessel}[1][mug]{the #1}
\end{codeblocktitled}

A command for a symbol you use often, such as \code|\Tenv|, has two advantages:
less typing, and the symbol can be changed in one place.
The same works for environments, with \code|\newenvironment{name}{before}{after}|.

\subsection{The \texttt{@} in command names}

Inside packages and classes, internal commands have an \code|@| in their name, e.g.\ \code|\@title|,
so that they cannot clash with commands in a document.
In a document, \code|@| is not a letter, so such a name cannot be typed directly.
To change an internal command in the preamble, put the change between \code|\makeatletter| and \code|\makeatother|:

\begin{codeblocktitled}{Preamble}
\makeatletter
\renewcommand{\@seccntformat}[1]{\csname the#1\endcsname.\quad}
\makeatother
\end{codeblocktitled}

(This one puts a full stop after section numbers: \enquote{1.~Introduction}.)
In a \texttt{.sty} file (\autoref{sec:sty}), \code|@| is a letter anyway.

\subsection{Beyond \texttt{\textbackslash newcommand}}

Since 2020, \hologo{LaTeX} also has \code|\NewDocumentCommand|,
which can take several optional arguments, starred forms and more.
The example above would be:

\begin{codeblocktitled}{Preamble}
\NewDocumentCommand{\vessel}{O{mug}}{the #1}
\end{codeblocktitled}

Underneath it is the L3 programming layer, a programming language for package authors.
You do not need it for a report;
if you are curious, \code|texdoc usrguide| describes the commands for documents,
and \code|texdoc interface3| the programming layer.

\section{One file per chapter}\label{sec:input}

A larger document keeps one file per chapter in a folder,
and the main file only collects them (\autoref{fig:tree}):

\begin{figure}[htbp]
    \centering
    \begin{minipage}[c]{.36\linewidth}
        \centering
        \begin{forest}
            for tree={%
                folder,
                grow'=0,
                font=\ttfamily\small,
                inner xsep=1pt,
                inner ysep=1.5pt,
                s sep=1pt,
                edge={black!40, line width=.6pt},
                folder indent=.5em,
            },
            where n children=0{}{text=figaccent},
            [thesis/
                [main.tex]
                [preamble.tex]
                [references.bib]
                [figures/
                    [cooling-curves.pdf]
                ]
                [sections/
                    [introduction.tex]
                    [theory.tex]
                    [results.tex]
                ]
            ]
        \end{forest}
    \end{minipage}\hfill
    \begin{minipage}[c]{.6\linewidth}
        \tcbinputlisting{codebox, grow to left by=0pt, grow to right by=0pt,
            enlarge left by=0pt, enlarge right by=0pt, width=\linewidth, title=\texttt{main.tex},
            listing file={\slidesdir/listings/99-input.tex}, listing options={style=code, numbers=none}}
    \end{minipage}
    \caption{A document in several files.}
    \label{fig:tree}
\end{figure}

\code|\input{sections/theory.tex}| pastes the file in at that point, as if it were typed there.
A chapter file has no preamble,
neither \code|\documentclass| nor \code|\begin{document}|;
it starts right away with \code|\section{...}| (or \code|\chapter{...}|).
Paths start at the folder of the main file, as for images.
Only the main file is compiled; on Overleaf, it is the one set as \textbf{Main document} in the menu.

\subsection{\texttt{\textbackslash include} and \texttt{\textbackslash includeonly}}

\code|\include{sections/theory}| works like \code|\input|, with three differences:
it always starts a new page, it cannot be nested,
and the file name is written without \texttt{.tex}.
In return, \code|\includeonly{sections/theory}| in the preamble compiles only that chapter,
while page numbers and references to the other chapters stay correct.
For a long thesis with many images, that makes compiling while you write much faster.
For a report, \code|\input| is simpler.

\section{A separate preamble}\label{sec:preamble-file}

The preamble can move out of the main file the same way
(line 2 of \texttt{main.tex} in \autoref{fig:tree}).
\texttt{preamble.tex} contains only \code|\usepackage| lines, settings and your own commands;
\texttt{main.tex} keeps the document class.
Your next report then starts with a copy of \texttt{preamble.tex}:
the \enquote{copy your last preamble} advice of \autoref{chap:resources}, made easy.

\section{A package of your own}\label{sec:sty}

Some templates go one step further and put the preamble into a package of their own, a \texttt{.sty} file,
loaded like any other package:

\begin{codeblock}
\usepackage{mystyle}
\end{codeblock}

The file \texttt{mystyle.sty} sits next to the main file.
It is a preamble with a few conventions on top:

\codefiletitled{listings/ch10-mystyle.sty}{\texttt{mystyle.sty}}

\begin{itemize}
    \item \code|\NeedsTeXFormat| and \code|\ProvidesPackage| identify the file;
        the name must match the file name, and the date and description in brackets appear in the log.
    \item Other packages are loaded with \code|\RequirePackage| instead of \code|\usepackage|.
    \item \code|@| is a letter inside the file, without \code|\makeatletter|.
    \item \code|\endinput| ends the file; anything after it is ignored.
\end{itemize}

There is no real advantage over \code|\input{preamble.tex}| for a single document.
A \texttt{.sty} file becomes useful when several documents share it, e.g.\ the template of a group or a department,
or when it should take options, which is beyond this handout
(\code|texdoc clsguide| describes how).
